\begin{center}
  \section*{Abstract}
\end{center}

\addcontentsline{toc}{section}{Abstract}

\begin{spacing}{1.5}
The growth of Internet of Things (IoT) infrastructure in Saudi Arabia, especially under Vision 2030 smart-city initiatives, increases the need for practical systems that can observe and analyze attacks against exposed devices. Many IoT devices still rely on weak credentials, limited update mechanisms, and remote access services such as SSH and Telnet. These conditions make IoT environments attractive targets for botnets and automated attacks.

This project presents TrapPot, an AI-powered IoT honeypot system that combines Cowrie, Zeek, a Random Forest classifier, and the ELK Stack. Cowrie emulates SSH and Telnet services in an isolated Docker-based environment, while Zeek converts network traffic into structured connection logs. The AI detector uses twelve Zeek-compatible features from the trained model pipeline to classify traffic into five classes: brute-force, scanning, malware download, command execution, and normal traffic. Logstash, Elasticsearch, and Kibana support log processing, indexing, and dashboard-based monitoring.

The Random Forest model was trained using the IoT-23 dataset and achieved an accuracy of approximately 99.59\% on the prepared dataset. The implemented system was also tested locally through Docker Compose, confirming that Cowrie captured sessions, Zeek generated \texttt{conn.log} and \texttt{ssh.log}, the AI detector produced live classification output, and Kibana displayed the collected data through the TrapPot dashboard. TrapPot therefore provides a working proof of concept for IoT honeypot monitoring, network-based classification, and security visualization in a controlled academic environment.

  {\bf Keywords:} TrapPot, Internet of Things (IoT), Honeypot, Random Forest (RF), Cybersecurity, Artificial Intelligence (AI).
\end{spacing}
\clearpage
